\section{停留时间分布的统计特征值}

\begin{frame}
    \frametitle{停留时间分布的统计特征值}

    与其他统计分布一样，为了比较不同的停留时间分布，通常是比较其统计特征值。常用的统计特征值有两个。
    \begin{itemize}
        \item 数学期望
        
        数学期望也就是均值，对停留时间分布而言即平均停留时间$\bar t$
        \[
            \bar{t}=\mu_{1}=\frac{\int_{0}^{\infty} t E(t) \mathrm{d} t}{\int_{0}^{\infty} E(t) \mathrm{d} t}=\int_{0}^{\infty} t E(t) \mathrm{d} t
        \]
        \item 方差
        
        方差为对均值的二阶矩。停留时间分布的方差为
        \[
            \sigma_{\mathrm{t}}^{2}=\mu_{2}^{\prime}=\int_{0}^{\infty}(t-\bar{t})^{2} E(t) \mathrm{d} t=\int_{0}^{\infty} t^{2} E(t) \mathrm{d} t-\bar{t}^{2} 
        \]
        方差表示对均值的离散程度，方差越大，停留时间长短不一参差不齐的程度越大，即返混越大。
    \end{itemize}

\end{frame}

% 在概率论和统计学中，数学期望（或均值，亦简称期望）是试验中每次可能结果的概率乘以其结果的总和，是最基本的数学特征之一。它反映随机变量平均取值的大小。
% 设连续性随机变量X的概率密度函数为f(x)，若积分绝对收敛，则称积分$\int_{0}^{\infty} t E(t) \mathrm{d} t$的值为随机变量的数学期望


\begin{frame}
    \frametitle{停留时间分布的统计特征值}

    若采用无因次时间，可得无因次平均停留时间及无因次方差。
    \[
        \bar{\theta}=\int_{0}^{\infty} \theta E(\theta) \mathrm{d} \theta
    \]
    \[
        \sigma_{\theta}^{2}=\frac{\sigma_{\mathrm{t}}^{2}}{\bar{t}^{2}}=\int_{0}^{\infty} \theta^{2} E(\theta) \mathrm{d} \theta-1
    \]
    $\sigma_{\theta}^{2}$介于[0,1]之间。

    当$\sigma_{\theta}^{2} \to 0$时，返混最小；

    当$\sigma_{\theta}^{2} \to 1$时，返混最大。

\end{frame}


\begin{frame}[t]
    \frametitle{}

    \begin{example}
        用（1）脉冲法；（2）升阶法分别测得一流动系统的响应曲线$c(t)$，试推导平均停留时间$\bar{t}$及方差$\sigma_{t}^{2}$与$c(t)$的关系式。
    \end{example}
    \pause
    \begin{minipage}{0.32\linewidth}
        (1) 脉冲法
        \[
            \begin{aligned}
                \bar{t}&=\frac{\int_{0}^{\infty} t c(t) \mathrm{d} t}{\int_{0}^{\infty} c(t) \mathrm{d} t} \\
                \sigma_{\mathrm{t}}^{2}&=\frac{\int_{0}^{\infty} t^2 c(t) \mathrm{d} t}{\int_{0}^{\infty} c(t) \mathrm{d} t}-\bar{t}^{2} 
            \end{aligned}
        \]
    \end{minipage}\hfill
    \begin{minipage}{0.32\linewidth}
        (2) 升阶法
        \[
            \begin{aligned}
                \bar{t}&=\int_{0}^{T} \left[1-\frac{c(t)}{c(\infty)}\right] \mathrm{d} t \\
                \sigma_{\mathrm{t}}^{2}&=2\int_{0}^{T} \left[1-\frac{c(t)}{c(\infty)}\right] \mathrm{d} t-(\bar{t})^{2} 
            \end{aligned}
        \]
    \end{minipage}\hfill
    \begin{minipage}{0.32\linewidth}
        (3) 降阶法
        \[
            \begin{aligned}
                \bar{t}&=\int_{0}^{T} \frac{c(t)}{c(0)} \mathrm{d} t \\
                \sigma_{\mathrm{t}}^{2}&=2\int_{0}^{T} \frac{tc(t)}{c(0)} \mathrm{d} t-(\bar{t})^{2} 
            \end{aligned}
        \]
    \end{minipage}
\end{frame}


\begin{frame}[t]
    \frametitle{}
    \begin{example}
        用脉冲法测定一流动反应器的停留时间分布，得到出口流中示踪剂的浓度$c(t)$与时间$t$的关系如下
        \begin{table}
            \centering
            \begin{tabular}{cccccccccccccccc}
                \bottomrule
                $t$/min & 0 & 2 & 4 & 6 & 8 & 10 & 12 & 14 & 16 & 18 & 20 & 22 & 24 \\ \hline
                $c(t)$/(g/m$^3$) & 0 & 1 & 4 & 7 & 9 & 8 & 5 & 2 & 1.5 & 1 & 0.6 & 0.2 & 0 \\ \toprule
            \end{tabular}
        \end{table}
        试求平均停留时间及方差。
    \end{example}
\end{frame}